\chapterhead{2}{p-adic modular forms}

This chapter is devoted to the study of various properly p-adic generalizations
of the notion of modular form\index{modular form}, as "functions" of p-adic elliptic curves\index{elliptic curve} whose
\uline{Hasse invariant\index{Hasse invariant}} is not too near zero.

\sechead{2.0}{The Hasse invariant $A$ as a modular form; its q-expansion}

Let $R$ be any ring in which $p=0$ (i.e., $R$ is an $\F_p$-algebra) and consider
an elliptic curve\index{elliptic curve} $E/R$. The p'th power mapping $F_{\mathrm{abs}}$ is an
additive p-linear endomorphism of $\cO_E$, hence induces a p-linear
endomorphism of the R-module $H^1(E,\cO_E)$. If $\omega$ is a base of
$\w_{E/R}$, it determines the \uline{dual} base $\eta$ of $H^1(E,\cO_E)$, and we
define $A(E,\omega)\in R$ by setting
$F_{\mathrm{abs}}^*(\eta) = A(E,\omega)\cdot\eta$.

Replacing $\omega$ by $\lambda\omega$, $\lambda\in R^\times$ has the effect of
replacing $\eta$ by $\lambda^{-1}\eta$, and
$F_{\mathrm{abs}}^*(\lambda^{-1}\eta) =
\lambda^{-p}F_{\mathrm{abs}}^*(\eta) =
\lambda^{-p}\cdot A(E,\omega)\cdot\eta =
\lambda^{1-p}A(E,\omega)\cdot\lambda^{-1}\eta$, whence
$A(E,\lambda\omega) = \lambda^{1-p}\cdot A(E,\omega)$, which shows that
$A(E,\omega)$ is a modular form\index{modular form} of level one and weight p-1 defined over
$\F_p$. More intrinsically, we may interpret $F_{\mathrm{abs}}^*$ as an
R-linear homomorphism
$F_{\mathrm{abs}}^*\colon F_{\mathrm{abs}}^*(H^1(E,\cO_E)) =
(H^1(E,\cO_E))^{\otimes p}\to H^1(E,\cO_E)$, so as a \uline{section} of
$(\w_{E/R})^{\otimes p-1}$. In terms of the base $\omega$ of $\w$, this section
is $A(E,\omega)\cdot\omega^{\otimes p-1}$. To see that $A$ is holomorphic at
$\infty$, we simply note that the Tate curve\index{Tate curve} over $\F_p((q))$ is the
restriction of a plane curve $C$ over $\F_p[[q]]$, and that it canonical
differential $\omega_{can}$ is the restriction of a base over $\F_p[[q]]$ of
the dualizing sheaf of $C$. Thus $\omega_{can}$ determines the dual base
$\eta_{can}$ of $H^1(C,\cO_C)$ as $\F_p[[q]]$-module, and
$A(\mathrm{Tate}(q),\omega_{can})$ is just the \uline{matrix} of
$F_{\mathrm{abs}}^*$ on $H^1(C,\cO_C)$ with respect to the base $\eta_{can}$.
In particular, $A(\mathrm{Tate}(q),\omega_{can}\in\F_p[[q]]$.

An alternative method of establishing holomorphy is to use the fact that for any
elliptic curve $E/R$ over any base ring $R$, $H^1(E,\cO_E)$ is the tangent
space of $E/R$ at the origin, which is to say the R-module of all
translation-invariant \uline{derivations} of $E/R$, and that when $R$ is an
$\F_p$-algebra, the action of $F_{\mathrm{abs}}^*$ on $H^1(E,\cO_E)$ consists
of taking the p'th iterate of an invariant derivation. Now we use the fact that
there is a local parameter $t$ on the completion of the Tate curve along its
identity section in terms of which $\omega_{can} = dt/1+t$. Let $D$ be the
invariant derivation dual to $\omega_{can}$. Then $D(t)=1+t$, hence
$D(1+t)=1+t$, hence $D^n(1+t)=1+t$ for all $n\geq 1$. Over $\F_p$, $D^p$ is an
invariant derivation, and it agrees with $D$ on $\omega_{can}$, hence $D^p=D$,
hence $F_{\mathrm{abs}}^*(\eta_{can})=\eta_{can}$, and
$A(\mathrm{Tate}(q),\omega_{can})=1$.

\sechead{2.1}{Deligne's congruence $A\equiv E_{p-1}\bmod p$}

For any even integer $k\geq 4$, the Eisenstein series $E_k$ is the modular form\index{modular form}
over $\C$ of level one and weight $k$ whose q-expansion is

\[
1 - \frac{2k}{b_k}\sum \sigma_{k-1}(n)q^n, \quad
\sigma_{k-1}(n)=\sum_{\substack{d\mid n\\ d\geq 1}} d^{k-1}.
\]

As its q-expansion coefficients all lie in $\Q$, $E_k$ is defined over $\Q$
(by \eqref{ka:1.9.1}). For $k=p-1$, $p\geq 5$, the p-adic ordinal of
$\frac{-2(p-1)}{b_{p-1}}$ is 1, hence $E_{p-1}$ has q-expansion coefficients in
$\Q\cap\Z_p$. Thus it makes sense to reduce $E_{p-1}$ modulo $p$, obtaining a
modular form over $\F_p$, whose q-expansion is the constant 1. Hence
$A\equiv E_{p-1}\bmod p$, because both are modular forms of the same weight with
the same q-expansions.

For $p=2$ and 3, it is \uline{not possible} to lift $A$ to a modular form of
level one, holomorphic at $\infty$, over $\Q\cap\Z_p$. However, for $p=2$ and,
$3\leq n\leq 11$, $2\nmid n$ we may lift $A$ to a modular form of
\uline{level n} and weight 1, holomorphic at $\infty$, over $\Z[1/n]$
(by \eqref{ka:1.7.1}). For $p=3$ and any $n\geq 3$, $3\nmid n$ we may lift $A$ to a
modular form of \uline{level n} and weight 2, holomorphic at $\infty$, over
$\Z[1/n]$ (by \eqref{ka:1.7.1}).

For $p=2$ and $3\leq n\leq 11$, $n$ odd (resp. for $p=3$ and $n\geq 2$,
$3\nmid n$), we choose a modular form $E_{p-1}$ of weight $p-1$ and level $n$,
holomorphic at $\infty$, defined over $\Z[1/n]$, which lifts $A$.

\begin{kathmstar}{Remark.}
For p=2, there exists a lifting of $A$ to a modular form of level n over
$\Z[1/n]$ for n=3, 5, 7, 9, 11, and hence for any n divisible by one of
3, 5, 7, 11. But the author does not know whether $A$ lifts to a form of level n
for other n (even for n=13!). An \uline{alternative} approach to the
difficulties caused by p=2 and 3 \uline{might} be based on the observation that
the Eisenstein series $E_4 = 1 + 240\sum \sigma_3(n)q^n$ provides a level 1
lifting to $\Z$ of $A^4$ if p=2 (resp. of $A^2$ if p=3).
\end{kathmstar}

\sechead{2.2}{p-adic modular forms with growth conditions}

\kaitemb{2.2.0} Let $R_0$ be a p-adically complete ring (i.e.
$R_0\approx\varprojlim R_0/p^N R_0$), and choose an element $r\in R_0$. For any
integer $n\geq 1$, prime to $p$, (resp. $3\leq n\leq 11$ for p=2, and
$n\geq 2$ for p=2) we define the module $M(R_0,r,n,k)$ of p-adic modular forms\index{p-adic modular form}
over $R_0$ of growth r, level n and weight k: An element
$f\in M(R_0,r,n,k)$ is a rule which assigns to any triple $(E/S,\alpha_n,Y)$
consisting of:

\kaitem{2.2.1} an elliptic curve\index{elliptic curve} $E/S$, where $S$ is a $R_0$-scheme on which p
is nilpotent (i.e. $p^N=0$ for $N\gg 0$);

\kaitem{2.2.2} a level n structure $\alpha_n$;

\kaitem{2.2.3} a section $Y$ of $\w^{\otimes(1-p)}$ satisfying
$Y\cdot E_{p-1}=r$;

a section $f(E/S,\alpha_n,Y)$ of $(\w_{E/S})^{\otimes k}$ over $S$, which
depends only on the isomorphism class of the triple, and whose formation
commutes with arbitrary change of base of $R_0$-schemes $S'\to S$.

Equivalently, we may interpret $f$ as a rule which attaches to each quadruple
$(E/R,\omega,\alpha_n,Y)$ consisting of:

\kaitem{2.2.4} an elliptic curve $E/R$, $R$ an $R_0$-algebra in which p is
nilpotent;

\kaitem{2.2.5} a base $\omega$ of $\w_{E/R}$;

\kaitem{2.2.6} a level n-structure;

\kaitem{2.2.7} an element $Y\in R$ satisfying $Y\cdot E_{p-1}(E,\omega)=r$,

an element $f(E/R,\omega,\alpha_n,Y)$ in $R$, which depends only on isomorphism
class of the quadruple, whose formation commutes with extension of scalars of
V-algebras, and which satisfies the functional equation:

\kaitem{2.2.8} $f(E/R,\lambda\omega,\alpha_n,\lambda^{p-1}Y) =
\lambda^{-k}f(E/R,\omega,\alpha_n,Y)$ for $\lambda\in R^\times$.

By passage to the limit, we can allow $R$ to be a p-adically complete
$R_0$-algebra in the above definition.

\kaitem{2.2.9} We say that $f$ is holomorphic at $\infty$ if for each integer
$N\geq 1$, its value on
$(\mathrm{Tate}(q^n),\omega_{can},\alpha_n,
r(E_{p-1}(\mathrm{Tate}(q^n),\omega_{can}))^{-1})$, considered over
$\Z((q))\tens (R_0/p^N R_0)[\zeta_n]$ lies in
$\Z[[q]]\tens (R_0/p^N R_0)[\zeta_n]$, for each level $n$ structure $\alpha_n$.
We denote by $S(R_0,r,n,k)$ the submodule of $M(R_0,r,n,k)$ consisting of forms
holomorphic at $\infty$.

As formal consequence of the definitions, we have

\begin{kadispb}{2.2.10}
M(R_0,r,n,k) = \varprojlim M(R_0/p^N R_0,r,n,k).
\end{kadispb}

\begin{kadispb}{2.2.11}
S(R_0,r,n,k) = \varprojlim S(R_0/p^N R_0,r,n,k).
\end{kadispb}

\sechead{2.3}{Determination of $M(R_0,r,n,k)$ when $p$ is nilpotent in $R_0$}

\kaitemb{2.3.0} We begin by determining the universal triple
$(E/S,\alpha_n,Y)$ supposing that $p$ is nilpotent in $R_0$, and $n\geq 3$.
For notational convenience, let's denote $\w^{\otimes 1-p}$ by $\check{\cL}$. By
the definition of $M_n$, the functor
$\cF_{R_0,r,n}\colon S\longmapsto S$-isomorphism classes of triples
$(E/S,\alpha_n,Y)$ is the functor $\cF_{R_0,r,n}\colon S\longmapsto$
$\left\{\begin{aligned} &R_0\text{-morphisms } g\colon S\to M_n\tens R_0,\text{ together}\\
&\text{with a section } Y\text{ of } g^*(\check{\cL})\text{ verifying }
Y\cdot g^*(E_{p-1})=r \end{aligned}\right.$ which we may view as a sub-functor of the
functor $\cF_{R_0,n}\colon S\longmapsto$ $\{R_0\text{-morphisms } g\colon S\to M_n$,
plus a section $Y$ of $g^*(\check{\cL})\}$. This last functor is
\uline{representable}, by the $M_n\tens R_0$-scheme

\begin{center}
\begin{tikzcd}
\underline{\Spec}_{M_n\tens R_0}(\Sym(\check{\cL})) \arrow[d] \\
M_n\tens R_0
\end{tikzcd}
\end{center}

Indeed, we may cover $M_n\tens R_0$ by affine opens $\Spec(B_1)$ over which
$\check{\cL}$ admits an invertible section $\check{\ell}_1$, and cover $S$ by
affine opens $\Spec(A_{ij})$ such that $g|\Spec(A_{ij})$ factors through
$\Spec(B_1)$. Over $\Spec(B_1)$, $\underline{\Spec}(\Sym(\check{\cL}))$ is
$\Spec(B_1[\check{\ell}_1])$. A section $Y$ of $g^*(\check{\cL})$ determines an
element $Y\cdot g^*(\check{\ell}_1)$ of $A_{ij}$, and then a lifting of the
given homomorphism $g\colon B_1\to A_{ij}$ to a homomorphism
$\tilde{g}_{ij}\colon B_1[\check{\ell}_1]\to A_{ij}$ by the formula

\[
\tilde{g}_{ij}(\sum b_k(\check{\ell}_1)^k) =
\sum g(b_k)(Y\cdot g^*(\check{\ell}_1))^k.
\]

These $\tilde{g}_1$ piece together to define a morphism from $S$ to
$\underline{\Spec}(\Sym(\check{\cL}))$.

The subfunctor $\cF_{R_0,r,n}$ is then represented by the closed subscheme of
$\underline{\Spec}(\Sym(\check{\cL}))$ defined by the vanishing of $E_{p-1}-r$.
Thus the universal triple $(E/S,\alpha_n,Y)$ is just the inverse image on
$\underline{\Spec}(\Sym(\check{\cL}))$ of the universal elliptic curve\index{elliptic curve} with
level n structure over $M_n\tens R_0$, hence

\begin{kathm}{Proposition}{2.3.1.}
When p is nilpotent in $R_0$, and $n\geq 3$ is prime to p, there is a canonical
isomorphism
\[
\begin{aligned}
M(R_0,r,n,k) &=
H^0(\underline{\Spec}_{M_n\tens R_0}(\Sym(\check{\cL})/(E_{p-1}-r),
\w^{\otimes k})) \\
&= H^0(M_n\tens R_0,\bigoplus_{j\geq 0}
(\w)^{\otimes(k+j(p-1))}/(E_{p-1}-r)) \\
\text{(because } M_n \text{ is affine)} &= H^0(M_n\tens R_0,\bigoplus_{j\geq 0}
(\w)^{\otimes(k+j(p-r))}/(E_{p-1}-r) \\
&= \bigoplus_{j\geq 0} M(R_0,n,k+j(p-1))/(E_{p-1}-r).
\end{aligned}
\]
\end{kathm}

\sechead{2.4}{Determination of $S(R_0,r,n,k)$ when $p$ is nilpotent in $R_0$.}

\begin{kathm}{Proposition}{2.4.1.}
Let $n\geq 3$, $p\nmid n$. Under the isomorphism \eqref{ka:2.3.1}, the submodule
$S(R_0,r,n,k)\subset M(R_0,r,n,k)$ is the submodule
$H^0(\underline{\Spec}\,\Mbar_n\tens_{R_0}
(\Sym(\check{\cL})/(E_{p-1}-r)),\w^{\otimes k})$ of
$H^0(\underline{\Spec}\,M_n\tens_{R_0}
(\Sym(\check{\cL})/(E_{p-1}-r))$.
\end{kathm}

\begin{kaproof}
It suffices to treat the case in which $R_0\ni\zeta_n$. Then the ring of the
completion of $\Mbar_n\tens_{R_0}$ along $\infty$ is a finite number of copies
of $R_0[[q]]$, hence the ring of the completion of
$\Spec\,\Mbar_n\tens_{R_0}(\Sym(\check{\cL})/(E_{p-1}-r))$ along the inverse
image of $\infty$ is isomorphic to a finite number of copies of

\[
R_0[[q]]\cong R_0[[q]][Y]/(Y\cdot
E_{p-1}(\mathrm{Tate}(q^n),\omega_{can},\alpha_n)-r)
\]

(an \uline{isomorphism} because
$E_{p-1}(\mathrm{Tate}(q^n),\omega_{can},\alpha_n)$ is \uline{invertible} in
$R_0[[q]]$. Thus the condition that an element
$f\in H^0(\underline{\Spec}\,\Mbar_n\tens_{R_0}
(\Sym(\check{\cL})/(E_{p-1}-r)),\w^{\otimes k})$ have \uline{holomorphic}
q-expansions is precisely the condition that it extend to a section of
$\w^{\otimes k}$ over $\Spec\,\Mbar_n\tens_{R_0}
(\Sym(\check{\cL})/(E_{p-1}-r),\w^{\otimes k})$.
\end{kaproof}

\begin{kathmstar}{Remark 2.4.1.1.}
Analogously to \eqref{ka:2.3.1}, we have
\[
\begin{aligned}
&H^0(\underline{\Spec}\,\Mbar_n\tens_{R_0}
(\Sym(\check{\cL})/(E_{p-1}-r)),\w^{\otimes k}) \\
&= H^0(\Mbar_n\tens_{R_0},\w^{\otimes k}\tens
\Sym(\check{\cL})/(E_{p-1}-r)) \\
&= H^0(\Mbar_n\tens_{R_0},\bigoplus_j
\w^{k+j(p-1)}/(E_{p-1}-r)).
\end{aligned}
\]
\end{kathmstar}

\sechead{2.5}{Determination of $S(R_0,r,n,k)$ in the limit}

\begin{kathm}{Theorem}{2.5.1.}
Let $n\geq 3$, and suppose either that $k\geq 2$ or that $k=1$ and $n\leq 11$,
or that $k=0$ and $p\neq 2$, or that $k=0$, $p=2$, and $n\leq 11$. Let $R_0$ be
any p-adically complete ring ($R_0\xrightarrow{\sim}\varprojlim R_0/p^N R_0$),
and suppose $r\in R_0$ is not a zero divisor in $R_0$. Then the homomorphism
\end{kathm}

\begin{equation}\tag*{2.5.1.0}\label{ka:2.5.1.0}
\begin{tikzcd}
\varprojlim H^0(\Mbar_n,\bigoplus_{j\geq 0}\w^{k+j(p-1)})\tens_{\Z[1/n]}
(R_0/p^N R_0)/(E_{p-1}-r) \arrow[d] \\
S(R_0,r,n,k) = \varprojlim S(R_0/p^N R_0,r,n,k)
\end{tikzcd}
\end{equation}

is an isomorphism.

\begin{kaproof}
Let $\delta$ denote the quasicoherent sheaf $\bigoplus_{j\geq 0}
\w^{k+j(p-1)}$ on $\Mbar_n$, and put
$\delta_N=\delta\tens R_0/p^N R_0$. The inverse system of exact sequences

\begin{kadispb}{2.5.1.1}
0 \to \delta_N \xrightarrow{E_{p-1}-r} \delta_N \to
\delta_N/(E_{p-1}-r) \to 0
\end{kadispb}

gives an inverse system of six-term cohomology sequences

\begin{kadispb}{2.5.1.2}
0 \to H^0(\Mbar_n,\delta_N)\xrightarrow{E_{p-1}-r} H^0(\Mbar_n,\delta_N)
\to H^0(\Mbar_n,\delta_N/(E_{p-1}-r)) \to H^1(\Mbar_n,\delta_N) \\
\xrightarrow{E_{p-1}-r} H^1(\Mbar_n,\delta_N) \to
H^1(\Mbar_n,\delta_N/(E_{p-1}-r)) \to 0.
\end{kadispb}

Suppose first that $k>0$. Under our hypotheses, the base-changing theorem
\eqref{ka:1.7.1} applies, according to which
$H^0(\Mbar_n,\delta_N)=H^0(\Mbar_n,\delta)\tens(R_0/p^N R_0)$, and
$H^1(\Mbar_n,\delta_N)=0$. Thus the $H^0$ terms in \eqref{ka:2.5.1.2} form a short exact
sequence of inverse systems, the first of which has surjective transition
morphisms. Hence the inverse limits of these inverse systems form the desired
short exact sequence.

In case $k=0$ and $p\neq 2$ or $k=0$, $p=2$ and $n\leq 11$, we have
$H^1(\Mbar_n,\w^{\otimes k})=0$ for $k\geq 1$, hence
$H^1(\Mbar_n,\delta)=H^1(\Mbar_n,\cO)$, and by \eqref{ka:1.7.1},
$H^0(\Mbar_n,\delta_N)=H^0(\Mbar_n,\delta)\tens R_0/p^N R_0$. The exact
sequence \eqref{ka:2.5.1.2} becomes

\begin{kadispstar}
0 \to H^0(\Mbar_n,\delta)\tens R_0/p^N R_0 \to
H^0(\Mbar_n,\delta)\tens R_0/p^N R_0 \to \\
H^0(\Mbar_n,\delta_N/(E_{p-1}-r)) \to
H^1(\Mbar_n,\cO)\tens R_0/p^N R_0 \xrightarrow{r} \\
H^1(\Mbar_n,\cO)\tens R_0/p^N R_0 \to
H^0(\Mbar_n,\cO)\tens R_0/(p^N,r) \to 0.
\end{kadispstar}

For variable $N$, these form a six-term exact sequence of inverse systems. If
the sequence of their inverse limits were exact, the theorem would follow,
because the map
$\varprojlim H^0(\Mbar_n,\cO)\tens R_0/p^N R_0\xrightarrow{-r}
\varprojlim H^0(\Mbar_n,\cO)\tens R^0/p^N R_0$ is injective (this because
$H^0(\Mbar_n,\cO)$ is a finite free $\Z[1/n]$-module, and $r$ is not a zero
divisor in $R_0\xrightarrow{\sim}\varprojlim R_0/p^N R_0$). To prove the
exactness we apply a general lemma.
\end{kaproof}

\begin{kathm}{Lemma}{2.5.2.}
Let $0\to K^0\to K^1\to K^2\to\cdots$ be a (long) exact sequence in the
category of projective systems of abelian groups indexed by the positive
integers. Suppose that for all $i\neq i_0$, the projective system $K^i$ has
surjective transition morphisms, and that the sequence
$\varprojlim K^{i_0+1}\to\varprojlim K^{i_0+2}\to\varprojlim K^{i_0+3}$ is
exact. Then the sequence
$0\to\varprojlim K^0\to\varprojlim K^1\to\varprojlim K^2\to\cdots$ is exact.
\end{kathm}

\begin{kaproof}
Consider the 2 spectral sequences of hypercohomology for the functor
$\varprojlim$.

\[
\begin{aligned}
{}^{\mathrm{I}}E_2^{p,q} &=
H^p(R^q(\varprojlim)(K^\bullet)) \implies
\mathbb{R}^{p+q}(\varprojlim)(K^\bullet) \\
{}^{\mathrm{II}}E_2^{p,q} &=
R^p(\varprojlim)(H^q(K^\bullet)) \implies
\mathbb{R}^{p+q}(\varprojlim)(K^\bullet)
\end{aligned}
\]

By hypothesis, we have ${}^{\mathrm{II}}E_2^{p,q}=0$ for all values of $q$,
hence $\mathbb{R}^n(\varprojlim)(K^\bullet)=0$ for all $n$. According to
(\cite{roos61}), we have $R^i(\varprojlim)=0$ for $i\geq 2$, hence
${}^{\mathrm{I}}E_2^{p,q}=0$ for $q\geq 2$. By (\cite{roos61}), we have
$R^1(\varprojlim)(K^i)=0$ for $i\neq i_0$, hence

\[
{}^{\mathrm{I}}E_2^{p,q} = 0 \quad\text{unless}\quad q=0 \quad\text{or}\quad
q=1 \text{ and } p=i_0.
\]

As we have also supposed that ${}^{\mathrm{I}}E_2^{i_0+2,0}=0$, we have
degeneration: $E_2^{p,q}=E_\infty^{p,q}$ for all p,q. As $E_\infty^{p,q}=0$ for
all p,q, we get in particular ${}^{\mathrm{I}}E_2^{p,0}=0$ for all $p$, which
is the desired conclusion.
\end{kaproof}

\sechead{2.6}{Determination of a "basis" of $S(R_0,r,n,k)$ in the limit}

\begin{kathm}{Lemma}{2.6.1.}
Under the numerical hypotheses of theorem \eqref{ka:2.5.1}, for each $j\geq 0$ the
injective homomorphism

\begin{kadispb}{2.6.1.1}
H^0(\Mbar_n\tens\Z_p,\w^{\otimes k+j(p-1)})\xrightarrow{E_{p-1}}
H^0(\Mbar_n\tens\Z_p,\w^{\otimes k+(j+1)(p-1)})
\end{kadispb}

admits a section.
\end{kathm}

\begin{kaproof}
We must show that the cokernel of \eqref{ka:2.6.1.1} is a finite free $\Z_p$-module. By
the base-changing theorem \eqref{ka:1.7.1}, we have for each $j\geq 0$ an exact sequence
of finite free $\Z_p$-modules

\begin{kadispb}{2.6.1.1.1}
0 \to H^0(\Mbar_n\tens\Z_p,\w^{\otimes k+j(p-1)})\xrightarrow{E_{p-1}}
H^0(\Mbar_n\tens\Z_p,\w^{k+(j+1)(p-1)}) \to \\
\to H^0(\Mbar_n\tens\Z_p,\w^{\otimes k+(j+1)(p-1)}/
E_{p-1}\w^{\otimes k+j(p-1)}) \to
H^1(\Mbar_n\tens\Z_p,\w^{k+j(p-1)}) \to 0
\end{kadispb}

whose formation commutes with arbitrary change of base (for
$\w^{\otimes k+(j+1)(p-1)}\allowbreak/E_{p-1}\w^{\otimes k+j(p-1)}$, remark that it's
$\Z_p$-flat by Igusa's theorem\index{Igusa's theorem} (cf\cite{igusa58}), and modulo $p$, it becomes a skyscraper
sheaf on $M_n\tens\F_p$, hence has vanishing $H^1$). Hence the cokernel of the
map \eqref{ka:2.6.1.1} is the kernel of a surjective map of finite free $\Z_p$-modules,
hence is itself a finite free $\Z_p$-module.
\end{kaproof}

For each $n$, $k$ satisfying the hypotheses of \eqref{ka:2.5.1}, and each $j\geq 0$ we
choose once and for all a section of \eqref{ka:2.6.1.1}, and denote its image by
$B(n,k,j+1)$. Thus for $j\geq 0$, we have a direct sum decomposition

\begin{kadispb}{2.6.1.2}
H^0(\Mbar_n,\w^{\otimes k+(j+1)(p-1)})\simeq
E_{p-1}\cdot H^0(\Mbar_n,\w^{k+j(p-1)})\oplus B(n,k,j+1)
\end{kadispb}

and

\begin{kadispb}{2.6.1.3}
H^0(\Mbar_n,\w^{\otimes k})\xlongequal{\mathrm{dfn}} B(n,k,0)
\end{kadispb}

We define $B(R_0,n,k,j)=B(n,k,j)\tens R_0$. Iterating the $R_0$-analogue of
\eqref{ka:2.6.1.2} gives a direct sum decomposition

\begin{kadispb}[2.6.1.3-b]{2.6.1.3}
S(R_0,n,k+j(p-1))\xleftarrow{\sim}
\bigoplus_{a=0}^{j} B(R_0,n,k,a) \\ \sum E_{p-1}^{j-a}b_a \leftarrow \sum b_a
\end{kadispb}

Let $B^{\mathrm{rigid}}(R_0,r,n,k)$ denote the $R_0$-module consisting of all
formal sums

\[
\sum_{a=0}^{\infty} b_a, \quad b_a\in B(R,n,k,a)
\]

whose terms tend to zero in the sense that given any $N>0$, $\exists M>0$ such
that $b_a\in p^N\cdot B(R,n,k,a)$ for $a\geq M$, the M allowed to depend both
upon N and upon the series $\sum b_a$. (Notice that
$B^{\mathrm{rigid}}(R_0,r,n,k)$ does not depend upon r!)

\begin{kathm}{Proposition}{2.6.2.}
Hypotheses as in \eqref{ka:2.5.1}, the inclusion of $B^{\mathrm{rigid}}(R_0,r,n,k)$ in
the p-adic completion of $H^0(\Mbar_n,\bigoplus_{j\geq 0}
\w^{k+j(p-1)})$ induces (via \eqref{ka:2.6.1.3}) an \uline{isomorphism}

\begin{kadispb}{2.6.2.1}
B^{\mathrm{rigid}}(R_0,r,n,k) \longrightarrow S(R_0,r,n,k) \\
\sum b_a \longrightarrow \sum_{a\geq 0}\frac{r^a\cdot b_a}{E_{p-1}^a}
\end{kadispb}

where $\sum_{a\geq 0}\frac{r^a\cdot b_a}{E_{p-1}^a}$ has the value
$\sum_{a\geq 0} b_a(E/S,\alpha_n)\cdot Y^a$ on $(E/S,\alpha_n,Y)$.
\end{kathm}

\begin{kaproof}
For \uline{injectivity}, we must show that if
$\sum_{a\geq 0} b_a\in B^{\mathrm{rigid}}(R,n,k)$ can be written
$(E_{p-1}-r)\cdot\sum_{a\geq 0} s_a$ with
$s_a\in S(R,n,k+a(p-1))$, and $s_a$ tending to zero as $a\to\infty$, then all
$b_a=0$. It suffices to show that for any $N>0$, $b_a\equiv 0\bmod p^N$. But
$\bmod p^N$, both $\sum b_a$ and $\sum s_a$ become \uline{finite} sums. To fix
ideas, suppose $b_a\equiv s_a\equiv 0\bmod p^N$ $\forall a>M$. Let's show
$b_M\equiv s_M\equiv 0\bmod p^N$. As
$0\equiv b_{M+1}\equiv E_{p-1}s_M\bmod p^N$, $s_M\equiv 0\bmod p^N$, hence
$b_M\equiv E_{p-1}s_{M-1}\bmod p^N$, hence $b_M\equiv 0\bmod p^N$ by \eqref{ka:2.6.1.3}.
Now start again with M-1 ... .

For surjectivity, we just use the decomposition \eqref{ka:2.6.1.3}. Given $\sum s_a$,
$s_a\in S(R,n,k+a(p-1))$ tending to zero, we may decompose
$s_a=\sum_{i+j=a}(E_{p-1})^i b_j(a)$, with $b_j(a)\in B(R,n,k,j)$, and
$b_j(a)$ tends to zero as $a\to\infty$, \uline{uniformly} in j.

Then $\sum_a s_a = \sum_a\sum_{i+j=a}(E_{p-1})^i b_j(a) =
\sum_a\sum_{i+j=a} r^i b_j(a) + (E_{p-1}-r)\sum_a\sum_{i+j=a} b_j(a)
\sum_{u+v=i-1} E_{p-1}^u\cdot r^v$, hence $\sum s_a$ and
$\sum_a\sum_{i+j=a} r^i b_j(a)$ have the same image in $S(R_0,r,n,k)$. But for
each $j$, $\sum_i r^i b_j(i+j)$ converges to an element
$b_j^i\in B(R,n,k,j)$, and $b_j^i$ tends to zero as $j\to\infty$, and
$\sum_{j\geq 0} b_j^i$ has the same image in $S(R_0,r,n,k)$ as
$\sum_{a\geq 0} s_a$.
\end{kaproof}

\begin{kathm}{Corollary}{2.6.3.}
Hypotheses as in \eqref{ka:2.5.1}, the canonical mapping
$S(R_0,r,n,k)\longrightarrow S(R_0,1,n,k)$ defined modularly by composition with
the transformation of functors:
$(E/S,\alpha_n,Y)\longrightarrow(E/S,\alpha_n,rY)$, is injective; the
corresponding map
$B^{\mathrm{rigid}}(R_0,r,n,k)\longrightarrow B^{\mathrm{rigid}}(R_0,1,n,k)$ is
given by

\[
\sum b_a\longrightarrow\sum r^a b_a.
\]
\end{kathm}

\sechead{2.7}{Banach norm and q-expansion for $r=1$}

\begin{kathm}{Proposition}{2.7.1.}
Hypotheses as in \eqref{ka:2.5.1}, let $x\in R_0$ be any element which divides a power
$p^N$, $N\geq 1$, of p. Then the following conditions on an element
$f\in S(R_0,1,n,k)$ are equivalent, for $k\geq 0$:

(1) $f\in x\cdot S(R_0,1,n,k)$,

(2) the q-expansions of f all lie in $x\cdot R_0[\zeta_n][[q]]$,

(3) on each of the $\varphi(n)$ connected components of
$\Mbar_n\tens_{\Z[1/n]}\Z[1/n,\zeta_n]$, there is at least one cusp where the
q-expansion of f lies in $x\cdot R_0[\zeta_n][[q]]$.
\end{kathm}

\begin{kaproof}
Clearly (1) $\Longrightarrow$ (2) $\Longrightarrow$ (3). We will prove
(3) $\Longrightarrow$ (1). Because $r=1$, we have
$S(R_0/xR_0,1,n,k)\simeq B^{\mathrm{rigid}}(R_0/xR_0,1,n,k)\simeq
B^{\mathrm{rigid}}(R_0,1,n,k)/x\cdot B^{\mathrm{rigid}}(R_0,1,n,k)$,

so replacing $R_0$ by $R_0/xR_0$, we are reduced to the case $x=0$, and p
nilpotent in $R_0$. In that case $f\in B^{\mathrm{rigid}}(R_0,1,n,k)$ is a
finite sum $\sum_{a=0}^{M} b_a$, $b_a\in B(R_0,n,k,a)$, and it's q-expansion at
$(\mathrm{Tate}(q^n),\omega_{can},\alpha_n,(E_{p-1})^{-1})$ is that of

\[
\sum_{a=0}^{N} b_a\cdot(E_{p-1})^{-a} =
\frac{\sum_{a=0}^{N} b_a\cdot(E_{p-1})^{N-a}}{(E_{p-1})^N},
\]

hence by hypothesis, $\sum_{a=0}^{N} b_a(E_{p-1})^{N-a}$ has q-expansion
\uline{zero} at one or more cusps on each geometric connected component of
$\Mbar_n$, hence by the q-expansion principle\index{q-expansion principle} \eqref{ka:1.6.2},
$\sum_{a=0}^{N} b_a(E_{p-1})^{N-a}=0$. By \eqref{ka:2.6.1.3}, each $b_a=0$.
\end{kaproof}

\begin{kathm}{Proposition}{2.7.2.}
Let n, k, R satisfy the hypotheses of \eqref{ka:2.5.1}. Suppose given for each cusp
$\alpha$ of $\Mbar_n$ a power series $f_\alpha(q)\in R_0[\zeta_n][[q]]$. The
following conditions are equivalent:

1. The $f_\alpha$ are the q-expansions of an (necessarily unique) element
$f\in S(R_0,1,n,k)$.

2. For every power $p^N$ of p, there exists a positive integer
$M\equiv 0\bmod p^{N-1}$, and a "true" modular form\index{modular form}
$\varepsilon_N\in S(R_0,n,k+M(p-1))$ whose q-expansions are congruent
$\bmod p^N$ to the given $f_\alpha$.
\end{kathm}

\begin{kaproof}
(1) $\Longrightarrow$ (2). Replacing $R_0$ by $R_0/p^N R_0$, we may suppose p
nilpotent in $R_0$. We must show that the q-expansion of f is the q-expansions
of a true modular form of level $n$ and weight $k'\geq k$,
$k'\equiv k\bmod p^{N-1}(p-1)$. But as we saw above [cf\eqref{ka:2.7.1}], for $M\gg 0$,
and p nilpotent in $R_0$, f has the same q-expansions as
$g/(E_{p-1})^M$, g truly modular of weight $k+M(p-1)$. Multiplying top and
bottom by a suitable power of $E_{p-1}$, we may suppose
$M\equiv 0\bmod p^{N-1}$. Then the q-expansion congruence
$E_{p-1}(q)\equiv 1\bmod(p)$ at each cusp gives
$(E_{p-1})^{p^{N-1}}(q)\equiv 1\bmod(p^N)$, hence
$(E_{p-1})^M(q)\equiv 1\bmod(p^N)$, and hence $f\bmod p^N$ has the same
q-expansion as g.

(2) $\Longrightarrow$ (1). Multiplying necessary $\varepsilon_N$ by a power of
$(E_{p-1})^{p^{N-1}}$, we may assume that the weights $k+M_N(p-1)$ of the
$\varepsilon_N$ are \uline{increasing} with N. Let $\Delta_N=M_{N+1}-M_N$. Then
$(\varepsilon_{N+1}-\varepsilon_N\cdot(E_{p-1})^{\Delta_N})$ lies in
$p^N\cdot S(R_0,n,k+M_{N+1}(p-1))$ by the q-expansion principle \eqref{ka:1.6.2}, hence
$\sum_N(\varepsilon_{N+1}-\varepsilon_N\cdot(E_{p-1})^{\Delta_N})$
"converges" to an element of $S(R_0,1,n,k)$, whose q-expansions are congruent
modulo $p^N$ to those of $\varepsilon_N$.
\end{kaproof}

\sechead{2.8}{Bases for levels one and two}

Suppose $p\neq 2,3$. Then $E_{p-1}$ \uline{is} a modular form\index{modular form} of level one
which lifts the Hasse invariant\index{Hasse invariant}, and hence for any p-adically complete ring
$R_0\ni r$ and integer $n\geq 3$ prime to p, the group $\GL_2(\Z/n\Z)$ acts on
the functor $\cF_{R_0,r,n}$ [by $g(E/S,\alpha_n,Y)=(E/S,g\alpha_n,Y)$ on the set
$\cF_{R_0,r,n}(S)$], hence on $M(R_0,r,n,k)$ and on $S(R_0,r,n,k)$. Clearly
$M(R_0,r,1,k)$ is just the submodule $M(R_0,r,n,k)^{\GL_2(\Z/n\Z)}$ of
\uline{invariants} under this action, and $S(R_0,r,1,k)$ is the submodule
$S(R_0,r,n,k)^{\GL_2(\Z/n\Z)}$ of $S(R_0,r,n,k)$. Now suppose $n=3$ or $n=4$.
This choice has the advantage that $\GL_2(\Z/n\Z)$ then has order prime to p
(because $p\neq 2,3$), and $P=\frac{1}{\#\GL_2(\Z/n\Z)}\sum g$ is then a
projection onto the invariants. Using P we may also make the chosen section of
\eqref{ka:2.6.1.1} invariant by $\GL_2(\Z/3\Z)$, and define
$B(1,k,j)=B(n,k,j)^{\GL_2(\Z/n\Z)}=P(B(n,k,j))$,
$B(R_0,1,k,j)=B(1,k,j)\tens R_0=B(R_0,n,k,j)^{\GL_2(\Z/n\Z)}$. Similarly, we
define $B^{\mathrm{rigid}}(R_0,r,1,k)=P(B^{\mathrm{rigid}}(R_0,r,n,k))=
(B^{\mathrm{rigid}}(R_0,r,n,k))^{\GL_2(\Z/n\Z)}$; it is the subspace of
$B^{\mathrm{rigid}}(R_0,r,n,k)$ consisting of the elements $\sum b_a$ each of
whose terms $b_a$ is invariant by $\GL_2(\Z/n\Z)$.

Applying the projector P to \eqref{ka:2.6.2} gives:

\begin{kathm}{Proposition}{2.8.1.}
Let $p\neq 2,3$, $R_0$ a p-adically complete ring and $r\in R_0$ not a
zero-divisor. Then for each $k\geq 0$, the canonical mapping

\begin{kadispb}{2.8.1.0}
B^{\mathrm{rigid}}(R_0,r,1,k) \longrightarrow S(R_0,r,1,k) \\
\sum b_a \longrightarrow "\sum \frac{r^a b_a}{(E_{p-1})^a}"
\end{kadispb}

is an isomorphism.
\end{kathm}

Now suppose $p\neq 2$, and consider level two. Let
$E_{p-1}\in S(\Z[\frac{1}{2}],2,p-1)$ a lifting of the Hasse invariant. Because
the subgroup $G_1$ has order prime to p,
$G_1=\mathrm{Kernel}\colon\GL(\Z/4\Z)\longrightarrow\GL(2,\Z/2\Z)$,
considerations similar to the above provide a projector
$P_1=\frac{1}{\#G_1}\sum G_1$ from level 4 to level 2. We have
$M(R_0,r,2,k)=M(R_0,r,4,k)^{G_1}=P_1(M(R_0,r,4,k))$,
$S(R_0,r,2,k)=S(R_0,r,4,k)^{G_1}=P_1(S(R_0,r,4,k))$,
$B^{\mathrm{rigid}}(R_0,r,2,k)=B^{\mathrm{rigid}}(R_0,r,4,k)^{G_1}$, the
subspace of $B^{\mathrm{rigid}}(R_0,r,4,k)$ of elements $\sum b_a$ with each
$b_a$ invariant by $G_1$. Applying $P_1$ to \eqref{ka:2.6.2} we get:

\begin{kathm}{Proposition}{2.8.2.}
Let $p\neq 2$, $R_0$ a p-adically complete ring and $r\in R_0$ not a
zero-divisor. For each $k\geq 0$, the canonical mapping

\begin{kadispb}{2.8.2.0}
B^{\mathrm{rigid}}(R_0,r,2,k) \longrightarrow S(R_0,r,2,k) \\
\sum b_a \longrightarrow "\sum \frac{r^a b_a}{(E_{p-1})^a}"
\end{kadispb}

is an isomorphism.
\end{kathm}

Applying the projectors $P$ or $P_1$ to \eqref{ka:2.7.1} gives

\begin{kathm}{Proposition}{2.8.3.}
Let $R_0$ be a p-adically complete ring. Suppose either that $p\neq 2$ and
$n=2$ or that $p\neq 2,3$ and $n=1$. Let $x\in R_0$ be any element which divides
a power $p^N$, $N\geq 1$ of p. The following conditions on an element
$f\in S(R_0,1,n,k)$ are equivalent:

(1) $f\in x\cdot S(R_0,1,n,k)$,

(2) the q-expansions of f all lie on $xR_0[[q]]$.
\end{kathm}

\sechead{2.9}{Interpretation via formal schemes}

Let $n\geq 3$, $p\nmid n$, $R_0$ a p-adically complete ring, and $r\in R_0$. We
denote by $M_n(R_0,r)$ (resp. $\widetilde{M}_n(R_0,r)$) the formal scheme over
$R_0$ given the compatible family of $R_0/p^N R_0$-schemes
$\underline{\Spec}_{M_n\tens R_0/p^N R_0}(\Sym(\check{\cL})/(E_{p-1}-r))$
(resp. $\underline{\Spec}_{\widetilde{M}_n\tens R_0/p^N R_0}
(\Sym(\check{\cL})/(E_{p-1}-r))$). We have

\begin{kadispstar}
M(R_0,r,n,k) = H^0(M_n(R_0,r),\w^{\otimes k}) \\
S(R_0,r,n,k) = H^0(\widetilde{M}_n(R_0,r),\w^{\otimes k}).
\end{kadispstar}

Equivalently, we may view $M_n(R_0,r)$ (resp. $\widetilde{M}_n(R_0,r)$) as the
completion along $p=0$ of the usual scheme
$\underline{\Spec}_{M_n\tens R_0}(\Sym(\check{\cL})/(E_p-r))$
(resp. $\underline{\Spec}_{\widetilde{M}_n\tens R_0}
(\Sym(\check{\cL})/(E_{p-1}-r))$. For any $r$, the first of these schemes is
affine, because $M_n$ is, and when r=1 \uline{both} schemes are affine.

The p-adic completions of their coordinate rings are just the rings
$M(R,r,n,0)$ and $S(R_0,1,n,0)$ respectively.
